\section{Two intersecting three-term progressions}\label{sec:double}

The support $F=\{0,1,2,4\}$ contains the progressions $\{0,1,2\}$ and
$\{0,2,4\}$. Both relations enter the weighted estimate. For a nonnegative
$x$ supported on $F$, put $w_i=x(i)^2$. The condition $x*x\leq 1$ is
equivalent to
\begin{equation}\label{d-feasibility}
\begin{gathered}
0\leq w_i\leq1,\qquad w_iw_j\leq\tfrac14\quad(i\ne j),\\
w_1+2\sqrt{w_0w_2}\leq1,\qquad
w_2+2\sqrt{w_0w_4}\leq1.
\end{gathered}
\end{equation}
All tuples in this section use the coordinate order $(0,1,2,4)$. We shall
bound nonnegative linear functionals of $w$ by their values on the profiles
\begin{equation}\label{d-profiles}
\begin{aligned}
H_0&=(\tfrac12,0,\tfrac12,\tfrac12),&
H_L&=(\tfrac12,\tfrac12,\tfrac18,\tfrac12),\\
H_R&=(\tfrac18,\tfrac12,\tfrac12,\tfrac12),&
V_L&=(1,\tfrac14,\tfrac9{64},\tfrac14),\\
V_R&=(\tfrac9{64},\tfrac14,1,\tfrac14),&
R&=(\tfrac9{64},\tfrac14,\tfrac14,1),
\end{aligned}
\end{equation}
and on feasible profiles with a zero coordinate. The tuples
$H_0,H_L,V_L,V_R$ are used only as upper bounds; they need not satisfy
\eqref{d-feasibility}.

\subsection{Profiles with a zero coordinate}

\begin{lemma}\label{d-boundary}
Suppose that $w$ satisfies \eqref{d-feasibility} and has a zero coordinate.
For every nonnegative kernel $k$ on $F$, put
$k_w=k\,1_{\operatorname{supp}(w)}$. Then
\begin{equation}\label{d-boundary-comparison}
(k\cdot w)^4\leq T_F(k_w)\leq T_F(k).
\end{equation}
The last comparison is strict when $k$ is positive on $F$. Moreover,
$\sum_iw_i\leq3/2$.
\end{lemma}

\begin{proof}
A three-term progression embeds in a four-term progression. Extending $k$
and $x$ by zero at the additional point preserves ordinary convolution
feasibility and the two quantities in the weighted comparison of
Section~\ref{sec:ap4}. This proves the first inequality when the remaining
support is a progression.

For a triangle with no three-term progression, its pair sums are distinct
apart from interchange. Write $a$ for its largest kernel value and $r$ for
the sum of the other two. Its weighted squared norm is
\begin{equation}\label{d-triangle-norm}
a+\frac r4.
\end{equation}
Indeed, if all squared coordinates are at most $1/2$, the half-triple
dominates them and is the average of the three tuples having one coordinate
equal to $1$ and the other two equal to $1/4$. Otherwise, a largest squared
coordinate $u\in[1/2,1]$ bounds the other two by $1/(4u)$. The resulting
linear functional is convex in $u$, so its maximum occurs at an endpoint.
Choosing the largest kernel coordinate for the unit coordinate proves
\eqref{d-triangle-norm}.

Fixing two copies of the point carrying $a$ in the fourfold maximum gives
six distinct sums and therefore
\begin{equation}\label{d-triangle-quartic}
T_F(k_w)
\geq a^4+a^3r+\frac23a^2r^2
\geq(a+r/4)^4.
\end{equation}
For the final inequality, expand the fourth power. Since $r\leq2a\leq3a$,
the quadratic and higher terms are at most
$(153/256)a^2r^2<(2/3)a^2r^2$ when $r>0$. The zero cases follow directly.
Supports of size at most two are covered by zero extension.

For a positive full kernel, a missing point $0,4,1$, or $2$ excludes the
sum $0,16,1$, or $14$, respectively, from the smaller fourfold sumset.
That sum has a positive contribution for the full kernel. This proves
strictness. Finally, the same pair-bound argument with all weights equal
to one gives total at most $3/2$ for any feasible three-point profile;
an additional progression constraint can only decrease this maximum.
\end{proof}

\subsection{Reduction of the feasible profiles}

An endpoint estimate will be used for maxima at the first and third
coordinates.

\begin{lemma}\label{d-endpoint-curves}
For $u\in[1/2,1]$, set
\begin{equation}\label{d-Jh}
J=\frac1{4u},\qquad h=J(1-J)^2.
\end{equation}
The curve $(u,J,h,J)$ is dominated coordinatewise by the convex hull of
$H_L,V_L,V_R$. The curve $(h,J,u,J)$ is dominated coordinatewise by the
convex hull of $H_R,V_R,V_L$.
\end{lemma}

\begin{proof}
Set
\begin{equation}\label{d-curve-weights}
\begin{gathered}
t=2u-1,\qquad s=\frac{(2u-1)(1-u)}{4u},\qquad \beta=\frac s2,\\
\alpha=t+\frac{23}{32}\beta,\qquad \gamma=1-\alpha-\beta.
\end{gathered}
\end{equation}
These weights are nonnegative and sum to one, since
\[
\gamma=(1-u)\left(2-\frac{55(2u-1)}{256u}\right)\geq0.
\]
In $\gamma H_L+\alpha V_L+\beta V_R$, the first coordinate equals $u$,
the second and fourth equal $J+201s/256$, and the third equals
\[
\frac{7+2u}{64}+\frac{1815s}{4096}.
\]
On the other hand,
\begin{equation}\label{d-curve-remainder}
h-\frac{7+2u}{64}
=s\frac{u^2+5u-1}{16u^2}\leq\frac{7s}{16},
\end{equation}
because $6u^2-5u+1=(3u-1)(2u-1)\geq0$. The inequality $1815>1792$
proves domination of the third coordinate. Interchanging the first and
third coordinates proves the second assertion.
\end{proof}

\begin{lemma}\label{d-reduction}
For every nonnegative $k$, its linear functional on the profiles satisfying
\eqref{d-feasibility} is bounded above by its largest value on the six
profiles in \eqref{d-profiles} and on the feasible profiles having a zero
coordinate.
\end{lemma}

\begin{proof}
First suppose all coordinates are at most $1/2$. Drop the second
progression constraint for an upper bound and increase $w_4$ to $1/2$.
For fixed $m=w_1\in[0,1/2]$, the first constraint gives
$w_0w_2\leq(1-m)^2/4$. Maximizing a nonnegative linear functional under
this product bound and the upper bounds $w_0,w_2\leq1/2$ reduces to
\[
(w_0,w_2)=\left(\frac12,\frac{(1-m)^2}{2}\right)
\quad\hbox{or its interchange}.
\]
To see this, on the product boundary the objective has the convex form
$a\ell+cK/\ell$; it is largest at an endpoint of the allowed interval.
The remaining objective is convex in $m$, so its endpoints give
$H_0,H_L,H_R$.

Otherwise choose a largest coordinate $u\geq1/2$. Each other coordinate
is at most $J=1/(4u)$. We consider its four possible locations.

\paragraph{Largest coordinate at index 1.}
Drop the second progression constraint and set $w_4=J$. The same
product-bound argument gives
\begin{equation}\label{d-middle-one}
(w_0,w_2)=(J,f)\quad\hbox{or its interchange},
\qquad f=u(1-u)^2.
\end{equation}
Assign the larger of $k(0),k(2)$ to $J$. The objective becomes
\[
k(1)u+k_{\rm small}(J+f)
+(k_{\rm large}-k_{\rm small}+k(4))J.
\]
It is convex in $u$, since $J''>0$ and
\[
(J+f)''=\frac1{2u^3}+6u-4>0.
\]
On $[1/2,2/3]$, the first summand is at least $27/16$ and the second
at least $-1$; above $2/3$ both are positive. The endpoint profiles are
$H_L,H_R$ and the feasible tuples
\[
(1/4,1,0,1/4),\qquad(0,1,1/4,1/4).
\]

\paragraph{Largest coordinate at index 0.}
Retain both progression constraints. For fixed $c=w_2\in[0,J]$, increase
the other two coordinates to
\begin{equation}\label{d-largest-zero-bounds}
w_1=\min\{J,1-\sqrt{c/J}\},\qquad w_4=J(1-c)^2.
\end{equation}
The first formula changes at $c=h$. The objective is convex on each of
$[0,h]$ and $[h,J]$: its nonaffine terms are a nonnegative quadratic
and, on the second interval, a nonnegative multiple of $-\sqrt c$.
The resulting endpoints are
\begin{equation}\label{d-largest-zero-profiles}
(u,J,0,J),\qquad
(u,J,h,J(1-h)^2),\qquad
(u,0,J,h).
\end{equation}
All three satisfy \eqref{d-feasibility}. The first and third have a zero
coordinate. The middle is dominated by $(u,J,h,J)$, to which
Lemma~\ref{d-endpoint-curves} applies.

\paragraph{Largest coordinate at index 2.}
Use $f=u(1-u)^2$ again. We have $f\leq h$: after multiplication by
$4u$, the equivalent inequality $1-J-2u(1-u)\geq0$ becomes
\[
(2u-1)(4u^2-2u+1)\geq0.
\]
For fixed $a=w_0>0$, the other admissible maxima are
\begin{equation}\label{d-largest-two-bounds}
w_1=\min\{J,1-2\sqrt{au}\},\qquad
w_4=\min\left\{J,\frac{(1-u)^2}{4a}\right\}.
\end{equation}
The objective is affine or convex on the intervals with endpoints
$0,f,h,J$, since the nonaffine terms are nonnegative multiples of
$1/a$ and $-\sqrt a$. On $[0,f]$ it increases with $a$. It is enough
to use the feasible profiles
\begin{equation}\label{d-largest-two-profiles}
(f,J,u,J),\qquad
\left(h,J,u,\frac{(1-u)^2}{4h}\right),\qquad
(J,0,u,f).
\end{equation}
When $u=1$, the condition $w_0w_4=0$ gives these same limiting endpoints
directly on its two branches. The last profile has a zero coordinate.
Since $h\geq f$, the middle one is dominated by $(h,J,u,J)$ and hence
by Lemma~\ref{d-endpoint-curves}. The first is dominated by the chord
from $H_R$ to the feasible tuple $(0,1/4,1,1/4)$, with weight $2u-1$
on the latter. Its first coordinate is $(1-u)/4\geq u(1-u)^2$;
its second and fourth dominate $J$ by convexity of $1/u$.

\paragraph{Largest coordinate at index 4.}
The second progression constraint gives
$w_0\leq J(1-w_2)^2$. Drop the first constraint for an upper bound.
The convex quadratic lies below its endpoint chord, so the profile is
dominated by a convex combination of
\begin{equation}\label{d-largest-four-curves}
B(u)=(J,J,0,u),\qquad W(u)=(J(1-J)^2,J,J,u).
\end{equation}
The weight on $W(u)$ is $w_2/J$. Both endpoints also satisfy the first
progression constraint; for $W(u)$ its left side is $3J-2J^2\leq1$
for $J\in[1/4,1/2]$.

Convexity of $1/u$ bounds $B(u)$ by the chord joining the feasible
tuples
\[
Q=(\tfrac12,\tfrac12,0,\tfrac12),\qquad
S=(\tfrac14,\tfrac14,0,1).
\]
For $W(u)$, put $Z=(1,0,1/4,9/64)$ and define
\begin{equation}\label{d-outer-weights}
\lambda=\frac{(1-u)(64u+55)}{87u},\qquad
\mu=\frac{(2u-1)(64u+23)}{87u},\qquad
\nu=\frac{32(2u-1)(1-u)}{87u}.
\end{equation}
The weights are nonnegative and sum to one. The tuple
$\lambda H_R+\mu R+\nu Z$ has coordinates $w_1=J$, $w_4=u$,
and $w_2=J+\nu/4$. Its first coordinate exceeds $J(1-J)^2$ by
\begin{equation}\label{d-outer-remainder}
\frac{(1-u)(2u-1)(576u^2-145u+29)}{1856u^3}\geq0.
\end{equation}
The quadratic is positive because its leading coefficient is positive
and its discriminant is $21025-66816<0$. The profiles $Z,Q,S$ are
feasible and have a zero coordinate. This finishes all four cases,
including ties.
\end{proof}

\subsection{Quartic comparisons}

Write $k=(A,B,C,D)$ on $(0,1,2,4)$. In the expansion of $(k\cdot v)^4$,
group monomials according to the sum of their four support indices.
Each group is at most its coefficient sum times the maximum defining
$T_F(k)$ at that index. Thus coefficient sums at most one prove the
desired comparison. The following computations use weighted AM--GM
to move a term to other fourfold products when a coefficient sum
exceeds one.

\begin{lemma}\label{d-half-comparisons}
For $v=H_0,H_L,H_R,R$ and $k\geq0$,
\begin{equation}\label{d-half-inequality}
(k\cdot v)^4\leq T_F(k).
\end{equation}
Each inequality is strict when $k$ is positive on $F$.
\end{lemma}

\begin{proof}
For $H_0$, the target is $(A+C+D)^4/16$ on the progression
$\{0,2,4\}$. After dividing indices by two, its coefficient sums are
\[
\frac1{16}(1,4,10,16,19,16,10,4,1).
\]
Use $A^2D^2\leq(A^3D+AD^3)/2$ on the entire term
$(6/16)A^2D^2$. The coefficient sums at indices $2,4,6$ become
$13/16$ each; all others were at most one. This compares with the
three-point maximum. The positive full kernel also has a contribution
at the missing odd sum $1$, giving strictness.

For $H_L$, its index polynomial is
$(1+z+z^2/4+z^4)/2$. If $c_n$ is the coefficient of its fourth power,
the complete coefficient sums are
\begin{equation}\label{d-HL-coefficients}
\begin{gathered}
\begin{array}{c rrrrrrrrr}
\toprule
n&0&1&2&3&4&5&6&7&8\\
\midrule
4096c_n&256&1024&1792&1792&2144&3520&3952&2576&2497\\
\bottomrule
\end{array}\\[6pt]
\begin{array}{c rrrrrrrr}
\toprule
n&9&10&11&12&13&14&15&16\\
\midrule
4096c_n&3264&2320&768&1120&1024&256&0&256\\
\bottomrule
\end{array}
\end{gathered}
\end{equation}
These follow by squaring the row
$(1,2,3/2,1/2,33/16,2,1/2,0,1)$ and dividing by $16$.
Every coefficient is at most $247/256<1$.

For $H_R$, the linear form is $A/8+(B+C+D)/2$. Its coefficient
sums $c_n$ are
\begin{equation}\label{d-HR-coefficients}
\begin{gathered}
\begin{array}{c rrrrrrrrr}
\toprule
n&0&1&2&3&4&5&6&7&8\\
\midrule
256c_n&1/16&1&7&28&71&124&172&224&262\\
\bottomrule
\end{array}\\[6pt]
\begin{array}{c rrrrrrrr}
\toprule
n&9&10&11&12&13&14&15&16\\
\midrule
256c_n&240&208&192&112&64&64&0&16\\
\bottomrule
\end{array}
\end{gathered}
\end{equation}
Only sum $8$ exceeds one. Replace its term $(3/128)A^2D^2$ by
$(3/256)A^3D+(3/256)AD^3$. The coefficient sums at indices $4,8,12$
become $74/256,1,115/256$, respectively. All others pass unchanged.
The coefficient $1/4096$ at index $0$ gives strictness when $A>0$.

For $R$, relabel the support by $p\mapsto4-p$. Its index polynomial becomes
\[
C(z)=1+\tfrac14z^2+\tfrac14z^3+\tfrac9{64}z^4.
\]
The coefficients of $C(z)^4$ at indices $0$ through $7$ are
\[
1,\quad0,\quad1,\quad1,\quad\tfrac{15}{16},\quad
\tfrac34,\quad\tfrac{55}{64},\quad\tfrac{39}{64}.
\]
For indices at least $8$, write $C=1+H$. In
$C^4=1+4H+6H^2+4H^3+H^4$, the first two terms do not contribute.
The contribution of $6H^2$ is at most $3/8$ and occurs only at index $8$.
Since $H$ is coefficientwise bounded by $z^2(1+z+z^2)/4$, the other
two contributions are at most $7/16$ and $19/256$. Here $7$ and $19$
are the largest coefficients of $(1+z+z^2)^3$ and $(1+z+z^2)^4$.
The total is at most $227/256<1$. Thus all coefficients pass, and the
coefficient $15/16$ at reflected index $4$ gives strictness for positive
kernels. Reflection reverses the output indices and preserves their sum.
\end{proof}

\begin{lemma}\label{d-VL-comparison}
For $k\geq0$, $(k\cdot V_L)^4\leq T_F(k)$, strictly when $k$ is positive
on $F$.
\end{lemma}

\begin{proof}
The linear form is $E=A+B/4+9C/64+D/4$. The coefficient sums in $E^4$
are listed below. An omitted entry is zero.
\begin{equation}\label{d-VL-ledger}
\begin{array}{r l r l}
\toprule
n&c_n&n&c_n\\
\midrule
0&1&9&3/16+243/16384\\
1&1&10&33/256+729/262144\\
2&15/16&11&27/1024\\
3&31/64&12&1/16+243/32768\\
4&1+467/2048&13&1/64\\
5&3/4+279/4096&14&9/1024\\
6&39/64+1215/65536&15&0\\
7&29/128+729/262144&16&1/256\\
8&3/8+351/4096+6561/16777216&&\\
\bottomrule
\end{array}
\end{equation}
Only sum $4$ exceeds one. Its polynomial is
\[
A^3D+\frac{243}{2048}A^2C^2+\frac{27}{256}AB^2C+\frac1{256}B^4
\leq U+aV+bW,
\]
where
\begin{equation}\label{d-VL-collision}
\begin{gathered}
U=A^3D,\quad V=A^2C^2,\quad W=B^4,\\
a=\frac{351}{2048},\qquad b=\frac{29}{512},\qquad
a+b=\frac{467}{2048}.
\end{gathered}
\end{equation}
This uses $AB^2C\leq(A^2C^2+B^4)/2$. For $M=\max(U,V,W)$,
\begin{equation}\label{d-minimum-split}
U+aV+bW\leq M+a\min(U,V)+b\min(U,W).
\end{equation}
Indeed,
$a(V-U)_++b(W-U)_+\leq(a+b)(M-U)\leq M-U$.
Weighted AM--GM also gives
\begin{equation}\label{d-scaled-minima}
\begin{aligned}
\min(A^3D,A^2C^2)&\leq\tfrac4{15}A^3C+\tfrac{25}{12}A^2D^2,\\
\min(A^3D,B^4)&\leq\tfrac4{15}A^2B^2+\tfrac{25}{12}A^2D^2.
\end{aligned}
\end{equation}
For the first line, the minimum is at most $U^{2/3}V^{1/3}$, which equals
$(A^3C)^{2/3}(A^2D^2)^{1/3}$. The stated bound follows since
$(4/15)^2(25/12)=4/27$. The second line has the same exponents.

The term $M$ is at most the maximum at sum $4$. All added terms from
\eqref{d-scaled-minima} go to sums $2$ and $8$. The final coefficient
sum at $2$ is
\begin{equation}\label{d-VL-two}
\frac{15}{16}+\frac{467}{2048}\frac4{15}
=\frac{7667}{7680}<1.
\end{equation}
The original coefficient sum at $8$ is less than $15/32$, and the added
weight is
\[
\frac{467}{2048}\frac{25}{12}=\frac{11675}{24576}<\frac12.
\]
Its final total is below $31/32$. The remaining sums pass by
\eqref{d-VL-ledger}. For positive kernels, the positive maximum and
remaining slack at sum $2$ give strictness.
\end{proof}

\begin{lemma}\label{d-VR-comparison}
For $k\geq0$, $(k\cdot V_R)^4\leq T_F(k)$, strictly when $k$ is positive
on $F$.
\end{lemma}

\begin{proof}
Relabel the support by $p\mapsto2-p$, and label the kernel values on positions
$(0,1,2,-2)$ by $(A,B,C,D)$. The profile is then
$(1,1/4,9/64,1/4)$, so the target is again $E=A+B/4+9C/64+D/4$.
Only sums $0,1,2,-2$ have coefficient sum greater than one.
Retain respectively
\begin{equation}\label{d-VR-retained}
A^4,\qquad A^3B,\qquad
\tfrac9{16}A^3C+\tfrac38A^2B^2,\qquad A^3D.
\end{equation}
The terms removed from these groups are
\begin{equation}\label{d-VR-removed}
\begin{array}{r l}
\toprule
\text{sum}&\text{removed terms}\\
\midrule
0&(3/16)AB^2D+(27/64)A^2CD+(243/32768)C^2D^2\\
1&(1/64)B^3D+(27/128)ABCD\\
2&(27/1024)B^2CD+(243/4096)AC^2D\\
-2&(3/128)B^2D^2+(27/256)ACD^2\\
\bottomrule
\end{array}
\end{equation}
Apply the following weighted AM--GM inequalities:
\begin{equation}\label{d-VR-redistribution}
\begin{aligned}
AB^2D&\leq(A^2D^2+B^4)/2,&
A^2CD&\leq(A^2C^2+A^2D^2)/2,\\
C^2D^2&\leq(C^4+D^4)/2,&
B^3D&\leq(3B^4+D^4)/4,\\
ABCD&\leq(A^2D^2+B^2C^2)/2,&
B^2CD&\leq B^4/2+C^4/4+D^4/4,\\
AC^2D&\leq(A^2D^2+C^4)/2,&
B^2D^2&\leq(B^4+D^4)/2,\\
ACD^2&\leq(A^2C^2+D^4)/2.&&
\end{aligned}
\end{equation}
Every destination is a fourfold product on the reflected support. Only
sums $4,-4,6,8,-8$ receive additions. Their coefficient sums are
\begin{equation}\label{d-VR-ledger}
\begin{array}{r l l l}
\toprule
\text{sum}&\text{original}&\text{addition}&\text{final}\\
\midrule
4&467/2048+729/262144&807/2048&163801/262144<5/8\\
-4&3/8+9/1024&3603/8192&6747/8192<7/8\\
6&1215/65536&27/256&8127/65536<1/8\\
8&6561/16777216&2619/65536&677025/16777216<1/16\\
-8&1/256&5155/65536&5411/65536<3/32\\
\bottomrule
\end{array}
\end{equation}
Besides these and the four groups in \eqref{d-VR-retained}, the
unchanged nonzero coefficient sums are
\begin{equation}\label{d-VR-unchanged}
\begin{array}{r l r l}
\toprule
\text{sum}&\text{coefficient sum}&\text{sum}&\text{coefficient sum}\\
\midrule
-6&1/16&3&31/64+243/16384\\
-5&1/64&5&279/4096\\
-3&3/16&7&729/262144\\
-1&3/4+27/1024&&\\
\bottomrule
\end{array}
\end{equation}
All are less than one. The retained groups at $0,1,-2$ have coefficient
sum one, and the group at $2$ has coefficient sum $15/16$. This proves
the inequality and gives strictness for positive kernels at sum $2$.
Reflection preserves the sum of the fourfold maxima.
\end{proof}

\begin{proposition}\label{d-certificate}
For $F=\{0,1,2,4\}$ and every nonnegative kernel $k$,
\begin{equation}\label{d-main-inequality}
N_F(k)^4\leq T_F(k).
\end{equation}
The inequality is strict if $k$ is positive on $F$.
\end{proposition}

\begin{proof}
Lemma~\ref{d-reduction} bounds the weighted objective by its values
on \eqref{d-profiles} and on feasible profiles with a zero coordinate.
Lemmas~\ref{d-boundary}, \ref{d-half-comparisons},
\ref{d-VL-comparison}, and~\ref{d-VR-comparison} bound the fourth power
of each of these values by $T_F(k)$.

For positive kernels every listed comparison is strict. The feasible
set is compact by \eqref{d-feasibility}, so its weighted maximum is
attained. Applying the strict comparison to a maximizing profile
proves the strict form of \eqref{d-main-inequality}.
\end{proof}

\subsection{The sharp unweighted norm}

The same reduction determines the unweighted norm exactly.

\begin{proposition}\label{d-sharp-norm}
For $F=\{0,1,2,4\}$,
\begin{equation}\label{d-norm-value}
\sup_{\substack{0\ne x\geq0\\ \operatorname{supp}(x)\subseteq F}}
\frac{\|x\|_2^2}{\|x*x\|_\infty}=\frac{105}{64}.
\end{equation}
Equality holds for $x=(3/8,1/2,1/2,1)$ in the coordinate order
$(0,1,2,4)$.
\end{proposition}

\begin{proof}
By homogeneity, normalize $\|x*x\|_\infty=1$. In
Lemma~\ref{d-reduction}, the totals of $V_L,V_R,R$ are $105/64$,
those of $H_L,H_R$ are $13/8=104/64$, and that of $H_0$ is $3/2$.
Every feasible profile with a zero coordinate has total at most $3/2$
by Lemma~\ref{d-boundary}. This gives the upper bound.

The squared profile of the stated $x$ is $R$. It satisfies
\eqref{d-feasibility}: the two progression left sides are $5/8$ and
$1$, all pair bounds hold, and the coordinate at $4$ is $1$.
Thus $\|x*x\|_\infty=1$ and $\|x\|_2^2=105/64$.
\end{proof}
